\section{Related Work}
\label{sec:related}

We organise related work by research objective.

\subsection{Semi data warehouse and STAC (O1)}

Data \emph{lakes} store raw heterogeneous payloads with schema-on-read flexibility, whereas data \emph{warehouses} enforce integrated schemas optimised for analytics.
We adopt a \textbf{semi data warehouse}: catalogue discipline and metadata completeness of a warehouse, while retaining heterogeneous ingest formats suited to UBEM (vectors, rasters, APIs, uploads).
The SpatioTemporal Asset Catalog (STAC)~\cite{stacspec} provides a JSON metadata model for geospatial assets that we extend with a custom \texttt{dl4r:} namespace in our ingest API.
Comparable catalogue approaches include Planetary Computer STAC, ISO~19115, and DCAT, but UBEM pipelines rarely expose stable \texttt{datasource\_id} values back-linked from feature calculators.

\subsection{UBEM ontology stacks (O2)}

Ontology-centric urban energy pipelines increasingly integrate BIM, GIS, and IoT.
Shi et al.~\cite{shi2023cim} propose OntoCIM with a five-step BIM--GIS--IoT integration workflow and SPARQL applications.
Wu et al.~\cite{wu2023bem} combine Brick and BOT for building energy modeling with rule-based thermal zoning.
Ma et al.~\cite{ma2024ubem} introduce Building Template Ontology (BTO) and Urban Building Ontology (UBO) for city-scale EnergyPlus workflows with GeoSPARQL.
Standards such as CityGML, 3DCityDB, BOT, SOSA/SSN, and GeoSPARQL provide domain vocabularies, yet prior work does not integrate operational CityDB schemas, calculator provenance, and LLM-maintained OBDA in one stack.

\subsection{Text-to-SQL benchmarks (O3)}

General benchmarks (Spider~\cite{yu2018spider}, BIRD~\cite{li2023bird}) lack PostGIS spatial functions and multi-schema CIM layouts (\texttt{cim\_vector}, \texttt{cim\_census}, \texttt{cim\_network}, \texttt{cim\_raster}).
Our \texttt{ai4db} pipeline generates taxonomy-labelled pairs covering 18 task types (spatial predicates, raster--vector joins, nested queries) and five domain types.

\subsection{Fine-tuned spatial-SQL LLMs (O4)}

Domain fine-tuning improves execution accuracy (EX) over frontier models on specialised SQL~\cite{survey_llm_t2sql_liu2025,wang_spatial_ijgi}.
Our MSc thesis (\emph{CIM Assist}) established SQLCoder-7B + QLoRA on a CIM PostGIS benchmark (84.6\% EX)~\cite{taherdoust2025cimassist,dettmers2023qlora}.
Paper~2 extends this capability from interactive database access to OBDA \emph{source}-SQL generation (O5).

\subsection{VKG automation and OBDA (O5)}

Virtual Knowledge Graphs (VKGs) expose RDF interfaces over relational data without full materialisation.
Bereta and Koubarakis~\cite{bereta2016ontop} extend Ontop for geospatial databases with GeoSPARQL-to-SQL translation.
LLM4VKG~\cite{xiao2025llm4vkg} automates VKG bootstrap using mapping patterns (SE, SR, SRm, SH) but targets generic schemas without domain-trained PostGIS SQL or provenance extensions.
GeoTriples~\cite{kyzirakos2018geotriples} emphasises practical geospatial mapping with user revision loops.

\subsection{Provenance and confidence in UBEM (O6)}

W3C PROV-O~\cite{prov-o} standardises provenance records.
UBEM uncertainty studies often report stock-level sensitivity but not per-building lineage tied to registered datasources.
Our contribution assigns confidence from source tier, cross-method agreement, input completeness, and datasource freshness---enabling auditability when ground truth is partial.

Table~\ref{tab:gap} synthesises gaps addressed by CIM Wizard 2.0.

\begin{table}[t]
  \centering
  \caption{Cross-objective research gap synthesis.}
  \label{tab:gap}
  \small
  \begin{tabular}{p{2.2cm}p{3.2cm}p{2.2cm}}
    \toprule
    Prior state & Limitation & Addressed by \\
    \midrule
    Ad hoc files/APIs & No datasource IDs & O1 \\
    Single ontology papers & No unified UBEM stack & O2 \\
    Generic text-to-SQL & No CIM PostGIS & O3 \\
    General LLMs & Low EX on spatial SQL & O4 \\
    Manual/generic VKG & No CIM provenance & O5 \\
    Priority-only execution & No confidence audit & O6 \\
    \bottomrule
  \end{tabular}
\end{table}

